% ============================================================
% Documento LaTeX - Site de Acolhimento FAESA
% Desenvolvimento Web II
% Compilar com: pdflatex ou lualatex (recomendado para TikZ)
% ============================================================
\documentclass[12pt, a4paper]{article}

% ---- Pacotes ----
\usepackage[utf8]{inputenc}
\usepackage[T1]{fontenc}
\usepackage[brazilian]{babel}
\usepackage{geometry}
\usepackage{graphicx}
\usepackage{xcolor}
\usepackage{hyperref}
\usepackage{tikz}
% \usepackage{tikz-uml}  % Removido: indisponível no Overleaf
\usepackage{enumitem}
\usepackage{booktabs}
\usepackage{longtable}
\usepackage{tabularx}
\usepackage{fancyhdr}
\usepackage{titlesec}
\usepackage{tcolorbox}
% \usepackage{pgf-umlsd}  % Removido: indisponível no Overleaf
\usepackage{float}
\usepackage{caption}
\usepackage{subcaption}
\usepackage{amsmath}
\usepackage{fontawesome5}
\usepackage{setspace}
\usepackage{parskip}

% ---- TikZ Libraries ----
\usetikzlibrary{shapes, shapes.geometric, shapes.multipart, arrows.meta, positioning, calc, fit, backgrounds, shadows, decorations.markings}

% ---- Cores FAESA ----
\definecolor{faesaAzul}{HTML}{003366}
\definecolor{faesaAzulClaro}{HTML}{0066CC}
\definecolor{faesaVerde}{HTML}{28A745}
\definecolor{faesaLaranja}{HTML}{FF8C00}
\definecolor{faesaCinza}{HTML}{6C757D}
\definecolor{faesaBranco}{HTML}{FFFFFF}
\definecolor{faesaBG}{HTML}{F5F7FA}
\definecolor{reqFunc}{HTML}{DCE6F1}
\definecolor{reqNFunc}{HTML}{FDE8D0}

% ---- Geometria ----
\geometry{
    top=2.5cm,
    bottom=2.5cm,
    left=2.5cm,
    right=2.5cm
}

% ---- Hyperref ----
\hypersetup{
    colorlinks=true,
    linkcolor=faesaAzul,
    urlcolor=faesaAzulClaro,
    citecolor=faesaAzul,
    pdfauthor={Gabriel Malheiros},
    pdftitle={Site de Acolhimento FAESA - Documento de Requisitos},
    pdfsubject={Desenvolvimento Web II}
}

% ---- Cabeçalho e Rodapé ----
\pagestyle{fancy}
\fancyhf{}
\fancyhead[L]{\textcolor{faesaAzul}{\small\textbf{FAESA — Site de Acolhimento}}}
\fancyhead[R]{\textcolor{faesaCinza}{\small Desenvolvimento Web II}}
\fancyfoot[C]{\thepage}
\renewcommand{\headrulewidth}{1pt}
\renewcommand{\headrule}{\hbox to\headwidth{\color{faesaAzul}\leaders\hrule height \headrulewidth\hfill}}

% ---- Títulos ----
\titleformat{\section}
    {\Large\bfseries\color{faesaAzul}}
    {\thesection.}{0.5em}{}[\color{faesaAzul}\titlerule]

\titleformat{\subsection}
    {\large\bfseries\color{faesaAzulClaro}}
    {\thesubsection}{0.5em}{}

\titleformat{\subsubsection}
    {\normalsize\bfseries\color{faesaCinza}}
    {\thesubsubsection}{0.5em}{}

% ---- Caixas coloridas ----
\newtcolorbox{destaque}[1][]{
    colback=faesaBG,
    colframe=faesaAzul,
    fonttitle=\bfseries,
    title=#1,
    rounded corners,
    boxrule=1pt,
    left=8pt, right=8pt, top=6pt, bottom=6pt
}

\newtcolorbox{alerta}[1][]{
    colback=reqNFunc,
    colframe=faesaLaranja,
    fonttitle=\bfseries,
    title=#1,
    rounded corners,
    boxrule=1pt,
    left=8pt, right=8pt, top=6pt, bottom=6pt
}

% ---- Estilos TikZ para diagramas ----
\tikzstyle{actor} = [draw, fill=faesaAzul!10, text=faesaAzul, font=\small\bfseries, minimum width=2cm, minimum height=0.8cm, rounded corners=3pt]
\tikzstyle{usecase} = [draw, fill=faesaBG, text=black, font=\small, ellipse, minimum width=3.5cm, minimum height=1.2cm, align=center, text width=3cm]
\tikzstyle{system} = [draw=faesaAzul, thick, dashed, rounded corners=10pt, inner sep=15pt]
\tikzstyle{component} = [draw, fill=faesaAzulClaro!15, text=black, font=\small\bfseries, minimum width=3cm, minimum height=1cm, rounded corners=5pt, drop shadow]
\tikzstyle{database} = [cylinder, draw, fill=faesaVerde!15, shape border rotate=90, aspect=0.25, minimum height=1.5cm, minimum width=2cm, font=\small\bfseries]
\tikzstyle{arrowstyle} = [-{Stealth[length=3mm]}, thick, color=faesaCinza]
\tikzstyle{dashedarrow} = [-{Stealth[length=3mm]}, thick, dashed, color=faesaCinza]
\tikzstyle{class} = [draw, fill=white, text=black, font=\small, rectangle split, rectangle split parts=3, minimum width=4cm, drop shadow]

% ============================================================
\begin{document}

% ---- Capa ----
\begin{titlepage}
\begin{center}

\vspace*{1cm}

% Logo simulado (substituir por imagem real)
\begin{tikzpicture}
    \node[fill=faesaAzul, rounded corners=10pt, minimum width=6cm, minimum height=2cm, text=white, font=\Huge\bfseries] {FAESA};
\end{tikzpicture}

\vspace{1.5cm}

{\LARGE\textcolor{faesaAzul}{\textbf{Centro Universitário FAESA}}}

\vspace{0.5cm}

{\large\textcolor{faesaCinza}{Curso de Sistemas de Informação / Ciência da Computação}}

\vspace{2cm}

\rule{\textwidth}{2pt}

\vspace{0.5cm}

{\Huge\textcolor{faesaAzul}{\textbf{Site de Acolhimento FAESA}}}

\vspace{0.3cm}

{\Large\textcolor{faesaAzulClaro}{Documento de Requisitos, Arquitetura e Diagramas}}

\vspace{0.5cm}

\rule{\textwidth}{2pt}

\vspace{2.5cm}

{\large\textbf{Disciplina:} Desenvolvimento Web II}

\vspace{0.3cm}

{\large\textbf{Aluno:} Gabriel Malheiros}

\vspace{0.3cm}

{\large\textbf{Professor(a):} [Nome do Professor]}

\vspace{\fill}

{\large Vitória --- ES}

{\large 2026}

\end{center}
\end{titlepage}

% ---- Sumário ----
\newpage
\tableofcontents
\newpage

% ============================================================
% SEÇÃO 1 - INTRODUÇÃO
% ============================================================
\section{Introdução}

\subsection{Contextualização}

A FAESA --- Centro Universitário, presente na educação capixaba há mais de 53 anos, é reconhecida pelo MEC como o melhor Centro Universitário do Espírito Santo pela sexta vez consecutiva e está entre os 10 melhores do Brasil (IGC/MEC). Com mais de 70 cursos de graduação e pós-graduação, mais de 75 mil alunos formados e 85\% dos egressos atuando em sua área de formação, a instituição é referência em ensino de qualidade.

Nesse contexto, o \textbf{programa de acolhimento estudantil} desempenha papel fundamental na integração dos novos alunos à vida acadêmica. O acolhimento abrange temas como:

\begin{itemize}[leftmargin=2cm]
    \item[\faIcon{book}] \textbf{Planos de Estudos} --- orientação para organização do tempo e metas acadêmicas;
    \item[\faIcon{chart-line}] \textbf{Melhoria nas Atividades} --- técnicas para aprimorar o desempenho em atividades curriculares e extracurriculares;
    \item[\faIcon{graduation-cap}] \textbf{Estudos fora da FAESA} --- incentivo a cursos complementares, certificações e aprendizado contínuo;
    \item[\faIcon{brain}] \textbf{Concentração e Foco} --- práticas de mindfulness, técnicas Pomodoro e exercícios de concentração;
    \item[\faIcon{users}] \textbf{Apoio Social e Emocional} --- conexão com colegas, mentores e suporte psicopedagógico.
\end{itemize}

\subsection{Objetivo do Documento}

Este documento tem como objetivo apresentar, de forma detalhada e organizada, todos os aspectos necessários para o desenvolvimento do \textbf{Site de Acolhimento da FAESA}, incluindo:

\begin{enumerate}[leftmargin=2cm]
    \item Levantamento completo de requisitos funcionais e não funcionais;
    \item Diagramas UML (casos de uso, classes, sequência, componentes);
    \item Arquitetura do sistema;
    \item Tecnologias recomendadas;
    \item Protótipos de interface;
    \item Cronograma de desenvolvimento.
\end{enumerate}

\subsection{Escopo do Projeto}

O site será uma \textbf{plataforma web responsiva} que centraliza recursos de acolhimento para os estudantes da FAESA, oferecendo ferramentas interativas para planejamento de estudos, acompanhamento de progresso, exercícios de concentração e acesso a conteúdos de apoio acadêmico.

\begin{destaque}[Visão do Produto]
    \textit{``Criar uma plataforma digital de acolhimento que acompanhe o estudante FAESA desde o ingresso até a conclusão do curso, promovendo organização, bem-estar e excelência acadêmica.''}
\end{destaque}

% ============================================================
% SEÇÃO 2 - ANÁLISE DE REQUISITOS
% ============================================================
\newpage
\section{Análise de Requisitos}

\subsection{Requisitos Funcionais (RF)}

Os requisitos funcionais descrevem as funcionalidades que o sistema deve oferecer aos seus usuários.

\begin{longtable}{|c|p{4.5cm}|p{6.5cm}|c|}
\hline
\rowcolor{faesaAzul}
\textcolor{white}{\textbf{ID}} & \textcolor{white}{\textbf{Requisito}} & \textcolor{white}{\textbf{Descrição}} & \textcolor{white}{\textbf{Prior.}} \\
\hline
\endfirsthead

\hline
\rowcolor{faesaAzul}
\textcolor{white}{\textbf{ID}} & \textcolor{white}{\textbf{Requisito}} & \textcolor{white}{\textbf{Descrição}} & \textcolor{white}{\textbf{Prior.}} \\
\hline
\endhead

\rowcolor{reqFunc}
RF01 & Cadastro e Login & O sistema deve permitir cadastro com matrícula FAESA e login via e-mail institucional (SSO) & Alta \\
\hline

RF02 & Plano de Estudos Personalizado & O aluno deve criar, editar e gerenciar seu plano de estudos com metas semanais e mensais & Alta \\
\hline

\rowcolor{reqFunc}
RF03 & Cronograma Interativo & Calendário interativo com drag-and-drop para organização de horários de estudo & Alta \\
\hline

RF04 & Exercícios de Concentração & Módulo com técnicas de concentração (Pomodoro, mindfulness, exercícios respiratórios) & Alta \\
\hline

\rowcolor{reqFunc}
RF05 & Dashboard de Progresso & Painel visual com gráficos de progresso acadêmico, metas atingidas e estatísticas & Alta \\
\hline

RF06 & Biblioteca de Recursos & Repositório de materiais sobre técnicas de estudo, artigos, vídeos e podcasts & Média \\
\hline

\rowcolor{reqFunc}
RF07 & Trilhas de Aprendizagem & Roteiros guiados de estudo para cada curso e período & Média \\
\hline

RF08 & Fórum de Discussão & Espaço para troca de experiências entre estudantes e mentores & Média \\
\hline

\rowcolor{reqFunc}
RF09 & Sistema de Mentoria & Conexão entre veteranos (mentores) e calouros para orientação & Média \\
\hline

RF10 & Notificações e Lembretes & Alertas via push e e-mail sobre metas, prazos e atividades pendentes & Média \\
\hline

\rowcolor{reqFunc}
RF11 & Avaliação de Bem-estar & Questionários periódicos para autoavaliação emocional e acadêmica & Média \\
\hline

RF12 & Atividades Extracurriculares & Catálogo de cursos, workshops e eventos externos recomendados & Baixa \\
\hline

\rowcolor{reqFunc}
RF13 & Gamificação & Sistema de pontos, badges e rankings para motivação dos estudantes & Baixa \\
\hline

RF14 & Relatórios para Coordenação & Relatórios agregados (anônimos) sobre engajamento e bem-estar estudantil & Baixa \\
\hline

\rowcolor{reqFunc}
RF15 & Chat com Suporte & Canal de comunicação direta com o núcleo de apoio psicopedagógico & Baixa \\
\hline

\caption{Requisitos Funcionais do Site de Acolhimento}
\label{tab:requisitos-funcionais}
\end{longtable}

\subsection{Requisitos Não Funcionais (RNF)}

Os requisitos não funcionais definem atributos de qualidade que o sistema deve possuir.

\begin{longtable}{|c|p{3.5cm}|p{7.5cm}|c|}
\hline
\rowcolor{faesaLaranja}
\textcolor{white}{\textbf{ID}} & \textcolor{white}{\textbf{Categoria}} & \textcolor{white}{\textbf{Descrição}} & \textcolor{white}{\textbf{Prior.}} \\
\hline
\endfirsthead

\hline
\rowcolor{faesaLaranja}
\textcolor{white}{\textbf{ID}} & \textcolor{white}{\textbf{Categoria}} & \textcolor{white}{\textbf{Descrição}} & \textcolor{white}{\textbf{Prior.}} \\
\hline
\endhead

\rowcolor{reqNFunc}
RNF01 & Responsividade & Interface adaptável a desktops, tablets e smartphones (Mobile-First) & Alta \\
\hline

RNF02 & Performance & Tempo de carregamento de páginas $\leq$ 3 segundos com conexão 3G & Alta \\
\hline

\rowcolor{reqNFunc}
RNF03 & Segurança & Autenticação OAuth 2.0 / SSO, criptografia TLS 1.3, proteção contra XSS e CSRF & Alta \\
\hline

RNF04 & Acessibilidade & Conformidade com WCAG 2.1 nível AA (contraste, leitores de tela, navegação por teclado) & Alta \\
\hline

\rowcolor{reqNFunc}
RNF05 & Disponibilidade & Uptime mínimo de 99,5\% com monitoramento 24/7 & Alta \\
\hline

RNF06 & Escalabilidade & Suportar crescimento de até 10.000 usuários simultâneos & Média \\
\hline

\rowcolor{reqNFunc}
RNF07 & Usabilidade & Interface intuitiva, seguindo Design System consistente, taxa de erro $\leq$ 2\% & Média \\
\hline

RNF08 & Manutenibilidade & Código documentado, arquitetura modular, cobertura de testes $\geq$ 80\% & Média \\
\hline

\rowcolor{reqNFunc}
RNF09 & LGPD & Conformidade total com a Lei Geral de Proteção de Dados (Lei 13.709/2018) & Alta \\
\hline

RNF10 & Internacionalização & Suporte a pt-BR como idioma principal, com preparação para en-US & Baixa \\
\hline

\caption{Requisitos Não Funcionais do Site de Acolhimento}
\label{tab:requisitos-nao-funcionais}
\end{longtable}

\subsection{Regras de Negócio}

\begin{alerta}[Regras de Negócio Essenciais]
\begin{enumerate}
    \item \textbf{RN01} --- Apenas estudantes com matrícula ativa na FAESA podem se cadastrar.
    \item \textbf{RN02} --- Os planos de estudo devem ser vinculados ao curso e período do aluno.
    \item \textbf{RN03} --- Mentores devem ser estudantes a partir do 5º período com CRA $\geq$ 7,0.
    \item \textbf{RN04} --- Os dados de bem-estar são confidenciais e acessados apenas de forma agregada pela coordenação.
    \item \textbf{RN05} --- O sistema de gamificação não pode gerar competição negativa; rankings são opcionais.
    \item \textbf{RN06} --- Conteúdos externos recomendados devem ser revisados por um coordenador antes da publicação.
\end{enumerate}
\end{alerta}

% ============================================================
% SEÇÃO 3 - PERSONAS E HISTÓRIAS DE USUÁRIO
% ============================================================
\newpage
\section{Personas e Histórias de Usuário}

\subsection{Personas}

\begin{destaque}[Persona 1 --- Lucas (Calouro)]
\begin{itemize}
    \item \textbf{Idade:} 18 anos
    \item \textbf{Curso:} Sistemas de Informação --- 1º período
    \item \textbf{Necessidade:} Organizar seus horários de estudo e se adaptar à rotina universitária
    \item \textbf{Dor:} Sente-se perdido com a quantidade de disciplinas e atividades
    \item \textbf{Motivação:} Quer ter um desempenho excelente desde o início
\end{itemize}
\end{destaque}

\begin{destaque}[Persona 2 --- Mariana (Veterana/Mentora)]
\begin{itemize}
    \item \textbf{Idade:} 22 anos
    \item \textbf{Curso:} Psicologia --- 7º período
    \item \textbf{Necessidade:} Compartilhar suas experiências e ajudar novos alunos
    \item \textbf{Dor:} Não há um canal estruturado para mentorias
    \item \textbf{Motivação:} Desenvolver habilidades de liderança e contribuir com a comunidade
\end{itemize}
\end{destaque}

\begin{destaque}[Persona 3 --- Prof. Ricardo (Coordenador)]
\begin{itemize}
    \item \textbf{Idade:} 45 anos
    \item \textbf{Função:} Coordenador de curso
    \item \textbf{Necessidade:} Acompanhar indicadores de engajamento e bem-estar dos alunos
    \item \textbf{Dor:} Dados dispersos e falta de visão consolidada
    \item \textbf{Motivação:} Reduzir evasão e melhorar a experiência acadêmica
\end{itemize}
\end{destaque}

\subsection{Histórias de Usuário}

\begin{center}
\begin{tabularx}{\textwidth}{|c|X|c|}
\hline
\rowcolor{faesaAzul}
\textcolor{white}{\textbf{ID}} & \textcolor{white}{\textbf{História de Usuário}} & \textcolor{white}{\textbf{RF}} \\
\hline
US01 & \textit{Como} calouro, \textit{quero} criar meu plano de estudos semanal \textit{para} me organizar melhor no semestre. & RF02 \\
\hline
US02 & \textit{Como} estudante, \textit{quero} praticar exercícios de concentração \textit{para} melhorar meu foco durante os estudos. & RF04 \\
\hline
US03 & \textit{Como} estudante, \textit{quero} visualizar meu progresso \textit{para} me manter motivado. & RF05 \\
\hline
US04 & \textit{Como} veterana, \textit{quero} me cadastrar como mentora \textit{para} ajudar calouros do meu curso. & RF09 \\
\hline
US05 & \textit{Como} coordenador, \textit{quero} visualizar relatórios de engajamento \textit{para} identificar alunos em risco. & RF14 \\
\hline
US06 & \textit{Como} estudante, \textit{quero} receber lembretes de metas \textit{para} não esquecer minhas tarefas. & RF10 \\
\hline
US07 & \textit{Como} estudante, \textit{quero} acessar uma biblioteca de recursos \textit{para} aprender técnicas de estudo eficientes. & RF06 \\
\hline
US08 & \textit{Como} estudante, \textit{quero} ganhar badges \textit{para} me sentir reconhecido pelo meu esforço. & RF13 \\
\hline
\end{tabularx}
\end{center}

% ============================================================
% SEÇÃO 4 - DIAGRAMAS UML
% ============================================================
\newpage
\section{Diagramas UML}

\subsection{Diagrama de Casos de Uso}

O diagrama de casos de uso apresenta as principais interações entre os atores do sistema (Estudante, Mentor, Coordenador e Administrador) e as funcionalidades oferecidas pela plataforma.

\begin{figure}[H]
\centering
\resizebox{\textwidth}{!}{%
\begin{tikzpicture}[node distance=1.2cm]

% Sistema boundary
\node[draw=faesaAzul, thick, dashed, rounded corners=15pt, 
      minimum width=13cm, minimum height=18cm, 
      label={[font=\Large\bfseries, text=faesaAzul]above:Sistema --- Site de Acolhimento FAESA}] 
      (sistema) at (6, -8) {};

% Atores à esquerda
\node[actor] (estudante) at (-3, -4) {\faIcon{user-graduate} Estudante};
\node[actor] (mentor) at (-3, -12) {\faIcon{chalkboard-teacher} Mentor};

% Atores à direita
\node[actor] (coord) at (15, -4) {\faIcon{user-tie} Coordenador};
\node[actor] (admin) at (15, -12) {\faIcon{user-cog} Admin};

% Casos de uso
\node[usecase] (uc1) at (4, -1) {Realizar Cadastro/Login};
\node[usecase] (uc2) at (8, -1) {Gerenciar Perfil};
\node[usecase] (uc3) at (4, -3.5) {Criar Plano de Estudos};
\node[usecase] (uc4) at (8, -3.5) {Gerenciar Cronograma};
\node[usecase] (uc5) at (4, -6) {Praticar Concentração};
\node[usecase] (uc6) at (8, -6) {Visualizar Dashboard};
\node[usecase] (uc7) at (4, -8.5) {Acessar Biblioteca};
\node[usecase] (uc8) at (8, -8.5) {Participar do Fórum};
\node[usecase] (uc9) at (4, -11) {Solicitar Mentoria};
\node[usecase] (uc10) at (8, -11) {Receber Notificações};
\node[usecase] (uc11) at (4, -13.5) {Responder Avaliação de Bem-estar};
\node[usecase] (uc12) at (8, -13.5) {Participar da Gamificação};
\node[usecase] (uc13) at (6, -16) {Visualizar Relatórios};

% Conexões Estudante
\draw[arrowstyle] (estudante) -- (uc1);
\draw[arrowstyle] (estudante) -- (uc3);
\draw[arrowstyle] (estudante) -- (uc5);
\draw[arrowstyle] (estudante) -- (uc7);
\draw[arrowstyle] (estudante) -- (uc9);
\draw[arrowstyle] (estudante) -- (uc11);
\draw[arrowstyle] (estudante) -- (uc6);
\draw[arrowstyle] (estudante) -- (uc10);
\draw[arrowstyle] (estudante) -- (uc12);

% Conexões Mentor
\draw[arrowstyle] (mentor) -- (uc1);
\draw[arrowstyle] (mentor) -- (uc8);
\draw[arrowstyle] (mentor) -- (uc9);

% Conexões Coordenador
\draw[arrowstyle] (coord) -- (uc1);
\draw[arrowstyle] (coord) -- (uc13);
\draw[arrowstyle] (coord) -- (uc6);

% Conexões Admin  
\draw[arrowstyle] (admin) -- (uc1);
\draw[arrowstyle] (admin) -- (uc2);
\draw[arrowstyle] (admin) -- (uc7);
\draw[arrowstyle] (admin) -- (uc13);

% Relacionamentos <<include>>
\draw[dashedarrow] (uc3) -- node[above, font=\tiny, sloped] {$\ll$include$\gg$} (uc4);
\draw[dashedarrow] (uc9) -- node[above, font=\tiny, sloped] {$\ll$include$\gg$} (uc8);

% Relacionamentos <<extend>>
\draw[dashedarrow] (uc12) -- node[above, font=\tiny, sloped] {$\ll$extend$\gg$} (uc6);
\draw[dashedarrow] (uc10) -- node[above, font=\tiny, sloped] {$\ll$extend$\gg$} (uc3);

\end{tikzpicture}
}%
\caption{Diagrama de Casos de Uso --- Site de Acolhimento FAESA}
\label{fig:casos-de-uso}
\end{figure}

% ============================================================
\newpage
\subsection{Diagrama de Classes}

O diagrama de classes modela as principais entidades do sistema e seus relacionamentos.

\begin{figure}[H]
\centering
\resizebox{\textwidth}{!}{%
\begin{tikzpicture}[
    every node/.style={font=\small},
    class/.style={draw, fill=white, rectangle split, rectangle split parts=3, 
                  minimum width=5cm, text centered, drop shadow, rounded corners=2pt},
    node distance=1.5cm
]

% Classe Usuario
\node[class] (usuario) at (0, 0) {
    \textbf{\textcolor{faesaAzul}{Usuario}}
    \nodepart{second}
    \begin{tabular}{l}
    - id: Long \\
    - nome: String \\
    - email: String \\
    - matricula: String \\
    - senha: String \\
    - tipo: TipoUsuario \\
    - dataCadastro: Date \\
    - ativo: Boolean
    \end{tabular}
    \nodepart{third}
    \begin{tabular}{l}
    + cadastrar(): void \\
    + autenticar(): Boolean \\
    + atualizarPerfil(): void
    \end{tabular}
};

% Classe PlanoEstudo
\node[class] (plano) at (8, 0) {
    \textbf{\textcolor{faesaAzul}{PlanoEstudo}}
    \nodepart{second}
    \begin{tabular}{l}
    - id: Long \\
    - titulo: String \\
    - descricao: String \\
    - dataInicio: Date \\
    - dataFim: Date \\
    - metaSemanal: Integer \\
    - status: StatusPlano
    \end{tabular}
    \nodepart{third}
    \begin{tabular}{l}
    + criar(): void \\
    + editar(): void \\
    + concluir(): void \\
    + calcularProgresso(): Double
    \end{tabular}
};

% Classe Meta
\node[class] (meta) at (16, 0) {
    \textbf{\textcolor{faesaAzul}{Meta}}
    \nodepart{second}
    \begin{tabular}{l}
    - id: Long \\
    - descricao: String \\
    - prazo: Date \\
    - concluida: Boolean \\
    - prioridade: Prioridade \\
    - categoria: String
    \end{tabular}
    \nodepart{third}
    \begin{tabular}{l}
    + marcarConcluida(): void \\
    + adiar(): void \\
    + notificar(): void
    \end{tabular}
};

% Classe SessaoConcentracao
\node[class] (sessao) at (0, -8) {
    \textbf{\textcolor{faesaAzul}{SessaoConcentracao}}
    \nodepart{second}
    \begin{tabular}{l}
    - id: Long \\
    - tipo: TipoExercicio \\
    - duracao: Integer \\
    - dataRealizacao: DateTime \\
    - pontuacao: Integer \\
    - concluida: Boolean
    \end{tabular}
    \nodepart{third}
    \begin{tabular}{l}
    + iniciar(): void \\
    + pausar(): void \\
    + finalizar(): void \\
    + registrarResultado(): void
    \end{tabular}
};

% Classe Mentoria
\node[class] (mentoria) at (8, -8) {
    \textbf{\textcolor{faesaAzul}{Mentoria}}
    \nodepart{second}
    \begin{tabular}{l}
    - id: Long \\
    - mentor: Usuario \\
    - mentorado: Usuario \\
    - dataInicio: Date \\
    - status: StatusMentoria \\
    - avaliacaoMentor: Integer \\
    - avaliacaoMentorado: Integer
    \end{tabular}
    \nodepart{third}
    \begin{tabular}{l}
    + solicitar(): void \\
    + aceitar(): void \\
    + encerrar(): void \\
    + avaliar(): void
    \end{tabular}
};

% Classe Recurso (Biblioteca)
\node[class] (recurso) at (16, -8) {
    \textbf{\textcolor{faesaAzul}{Recurso}}
    \nodepart{second}
    \begin{tabular}{l}
    - id: Long \\
    - titulo: String \\
    - tipo: TipoRecurso \\
    - url: String \\
    - descricao: String \\
    - tags: List$<$String$>$ \\
    - avaliacao: Double
    \end{tabular}
    \nodepart{third}
    \begin{tabular}{l}
    + publicar(): void \\
    + avaliar(): void \\
    + favoritar(): void
    \end{tabular}
};

% Classe AvaliacaoBemEstar
\node[class] (avaliacao) at (4, -15.5) {
    \textbf{\textcolor{faesaAzul}{AvaliacaoBemEstar}}
    \nodepart{second}
    \begin{tabular}{l}
    - id: Long \\
    - dataPreenchimento: Date \\
    - nivelEstresse: Integer \\
    - satisfacaoAcademica: Integer \\
    - horasSono: Double \\
    - respostas: Map$<$String, String$>$
    \end{tabular}
    \nodepart{third}
    \begin{tabular}{l}
    + preencher(): void \\
    + calcularIndice(): Double \\
    + gerarAlerta(): void
    \end{tabular}
};

% Classe Badge (Gamificação)
\node[class] (badge) at (12, -15.5) {
    \textbf{\textcolor{faesaAzul}{Badge}}
    \nodepart{second}
    \begin{tabular}{l}
    - id: Long \\
    - nome: String \\
    - descricao: String \\
    - icone: String \\
    - pontosNecessarios: Integer \\
    - categoria: CategoriaBadge
    \end{tabular}
    \nodepart{third}
    \begin{tabular}{l}
    + desbloquear(): void \\
    + verificarElegibilidade(): Bool
    \end{tabular}
};

% Relacionamentos
\draw[-{Stealth}, thick] (usuario) -- node[above, font=\scriptsize] {1..*} node[below, font=\scriptsize] {cria} (plano);
\draw[-{Stealth}, thick] (plano) -- node[above, font=\scriptsize] {1..*} node[below, font=\scriptsize] {contém} (meta);
\draw[-{Stealth}, thick] (usuario) -- node[left, font=\scriptsize] {0..*} node[right, font=\scriptsize] {realiza} (sessao);
\draw[-{Stealth}, thick] (usuario) -- node[right, font=\scriptsize, pos=0.3] {0..*} node[left, font=\scriptsize, pos=0.7] {participa} (mentoria);
\draw[-{Stealth}, thick] (usuario.south) -- ++(0,-1.5) -| node[above, font=\scriptsize, pos=0.25] {0..*} (avaliacao);
\draw[-{Stealth}, thick] (usuario.south) -- ++(0,-2) -| node[above, font=\scriptsize, pos=0.2] {0..*} (badge);
\draw[-{Diamond}, thick] (recurso) -- node[right, font=\scriptsize] {0..*} (meta);

\end{tikzpicture}
}%
\caption{Diagrama de Classes --- Entidades Principais do Sistema}
\label{fig:diagrama-classes}
\end{figure}

% ============================================================
\newpage
\subsection{Diagrama de Componentes / Arquitetura}

O diagrama a seguir ilustra a arquitetura de componentes do sistema, separando as responsabilidades em camadas.

\begin{figure}[H]
\centering
\resizebox{0.95\textwidth}{!}{%
\begin{tikzpicture}[node distance=1.5cm]

% ---- Camada de Apresentação ----
\node[draw=faesaAzul, fill=faesaAzul!5, thick, rounded corners=10pt, 
      minimum width=16cm, minimum height=3.5cm, 
      label={[font=\large\bfseries, text=faesaAzul]above left:Camada de Apresentação (Frontend)}] 
      (frontend) at (8, 0) {};

\node[component] (react) at (3, 0) {React.js / Next.js};
\node[component] (css) at (8, 0) {Tailwind CSS};
\node[component] (pwa) at (13, 0) {PWA / Service Worker};

% ---- Camada de API ----
\node[draw=faesaVerde, fill=faesaVerde!5, thick, rounded corners=10pt, 
      minimum width=16cm, minimum height=3.5cm, 
      label={[font=\large\bfseries, text=faesaVerde]above left:Camada de API (Backend)}] 
      (backend) at (8, -5) {};

\node[component, fill=faesaVerde!15] (api) at (3, -5) {API REST / GraphQL};
\node[component, fill=faesaVerde!15] (auth) at (8, -5) {Auth (JWT / OAuth)};
\node[component, fill=faesaVerde!15] (ws) at (13, -5) {WebSocket (Real-time)};

% ---- Camada de Serviços ----
\node[draw=faesaLaranja, fill=faesaLaranja!5, thick, rounded corners=10pt, 
      minimum width=16cm, minimum height=3.5cm, 
      label={[font=\large\bfseries, text=faesaLaranja]above left:Camada de Serviços (Domínio)}] 
      (servicos) at (8, -10) {};

\node[component, fill=faesaLaranja!15] (planoSvc) at (2, -10) {Plano de Estudos};
\node[component, fill=faesaLaranja!15] (concSvc) at (6.5, -10) {Concentração};
\node[component, fill=faesaLaranja!15] (mentSvc) at (10.5, -10) {Mentoria};
\node[component, fill=faesaLaranja!15] (gameSvc) at (14.5, -10) {Gamificação};

% ---- Camada de Dados ----
\node[draw=faesaCinza, fill=faesaCinza!5, thick, rounded corners=10pt, 
      minimum width=16cm, minimum height=3.5cm, 
      label={[font=\large\bfseries, text=faesaCinza]above left:Camada de Dados (Persistência)}] 
      (dados) at (8, -15) {};

\node[database] (postgres) at (3, -15) {PostgreSQL};
\node[database] (redis) at (8, -15) {Redis (Cache)};
\node[database] (s3) at (13, -15) {S3 / Storage};

% ---- Setas entre camadas ----
\draw[arrowstyle, very thick] (frontend) -- node[right, font=\small] {HTTP/HTTPS} (backend);
\draw[arrowstyle, very thick] (backend) -- node[right, font=\small] {Chamadas de Serviço} (servicos);
\draw[arrowstyle, very thick] (servicos) -- node[right, font=\small] {ORM / Queries} (dados);

% ---- Serviços Externos ----
\node[component, fill=red!10] (email) at (-2, -5) {E-mail Service};
\node[component, fill=red!10] (push) at (-2, -8) {Push Notifications};
\draw[dashedarrow] (backend) -- (email);
\draw[dashedarrow] (servicos) -- (push);

% Integrações
\node[component, fill=purple!10] (sso) at (18, -5) {FAESA SSO};
\node[component, fill=purple!10] (lms) at (18, -8) {LMS Integration};
\draw[dashedarrow] (backend) -- (sso);
\draw[dashedarrow] (servicos) -- (lms);

\end{tikzpicture}
}%
\caption{Diagrama de Componentes / Arquitetura do Sistema}
\label{fig:arquitetura}
\end{figure}

% ============================================================
\newpage
\subsection{Diagrama de Fluxo de Navegação}

Este diagrama apresenta o fluxo de navegação principal do usuário no site de acolhimento.

\begin{figure}[H]
\centering
\resizebox{0.9\textwidth}{!}{%
\begin{tikzpicture}[
    page/.style={draw=faesaAzul, fill=faesaAzul!10, rounded corners=8pt, 
                 minimum width=3.5cm, minimum height=1.8cm, text centered, 
                 font=\small\bfseries, text=faesaAzul, drop shadow},
    start/.style={draw=faesaVerde, fill=faesaVerde!20, circle, minimum size=1.5cm, 
                  font=\small\bfseries, text=faesaVerde},
    arrow/.style={-{Stealth[length=3mm]}, thick, color=faesaCinza}
]

% Início
\node[start] (inicio) at (0, 0) {Início};

% Páginas
\node[page] (login) at (4, 0) {Login / \\ Cadastro};
\node[page] (home) at (9, 0) {Home / \\ Dashboard};

% Ramificações a partir do Home
\node[page] (plano) at (4, -4) {Plano de \\ Estudos};
\node[page] (conc) at (9, -4) {Exercícios de \\ Concentração};
\node[page] (biblio) at (14, -4) {Biblioteca de \\ Recursos};

\node[page] (mentor) at (4, -8) {Mentoria};
\node[page] (forum) at (9, -8) {Fórum};
\node[page] (perfil) at (14, -8) {Meu Perfil};

\node[page] (bemestar) at (4, -12) {Avaliação de \\ Bem-estar};
\node[page] (gamif) at (9, -12) {Gamificação / \\ Badges};
\node[page] (ext) at (14, -12) {Atividades \\ Extracurriculares};

% Setas
\draw[arrow] (inicio) -- (login);
\draw[arrow] (login) -- (home);

\draw[arrow] (home) -- (plano);
\draw[arrow] (home) -- (conc);
\draw[arrow] (home) -- (biblio);

\draw[arrow] (home.south) -- ++(0, -0.5) -| (mentor);
\draw[arrow] (home.south) -- ++(0, -0.5) -| (forum);
\draw[arrow] (home.south) -- ++(0, -0.5) -| (perfil);

\draw[arrow] (plano) -- (mentor);
\draw[arrow] (conc) -- (forum);
\draw[arrow] (biblio) -- (perfil);

\draw[arrow] (mentor) -- (bemestar);
\draw[arrow] (forum) -- (gamif);
\draw[arrow] (perfil) -- (ext);

% Seta de volta
\draw[arrow, color=faesaVerde, dashed] (gamif.north) -- ++(0, 0.5) -- ++(5, 0) |- (home.east);

\end{tikzpicture}
}%
\caption{Diagrama de Fluxo de Navegação do Site}
\label{fig:fluxo-navegacao}
\end{figure}

% ============================================================
\newpage
\subsection{Diagrama Entidade-Relacionamento (ER)}

O diagrama ER apresenta o modelo de dados do sistema com suas entidades e relacionamentos.

\begin{figure}[H]
\centering
\resizebox{\textwidth}{!}{%
\begin{tikzpicture}[
    entity/.style={draw=faesaAzul, fill=faesaAzul!10, thick, 
                   minimum width=4cm, minimum height=1cm, rounded corners=5pt, 
                   font=\bfseries, text=faesaAzul, drop shadow},
    relation/.style={draw=faesaLaranja, fill=faesaLaranja!20, thick, 
                     diamond, minimum width=2.5cm, minimum height=1.2cm, 
                     font=\small\bfseries, text=faesaLaranja, aspect=2},
    attr/.style={draw=faesaCinza, fill=faesaCinza!10, ellipse, 
                 minimum width=2cm, font=\scriptsize, inner sep=2pt},
    pk/.style={draw=faesaCinza, fill=faesaCinza!10, ellipse, 
               minimum width=2cm, font=\scriptsize\underline, inner sep=2pt}
]

% Entidades
\node[entity] (user) at (0, 0) {USUARIO};
\node[entity] (plano) at (8, 3) {PLANO\_ESTUDO};
\node[entity] (meta) at (16, 3) {META};
\node[entity] (sessao) at (8, -3) {SESSAO\_CONCENTRACAO};
\node[entity] (mentoria) at (0, -6) {MENTORIA};
\node[entity] (recurso) at (16, -3) {RECURSO};
\node[entity] (avaliacao) at (8, -7) {AVALIACAO\_BEM\_ESTAR};
\node[entity] (badge) at (16, -7) {BADGE};

% Relacionamentos
\node[relation] (r1) at (4, 1.5) {cria};
\node[relation] (r2) at (12, 3) {contém};
\node[relation] (r3) at (4, -3) {realiza};
\node[relation] (r4) at (0, -3) {participa};
\node[relation] (r5) at (12, -5) {favorita};
\node[relation] (r6) at (4, -7) {preenche};
\node[relation] (r7) at (12, -7) {conquista};

% Conexões
\draw[thick] (user) -- node[above left, font=\scriptsize] {1} (r1);
\draw[thick] (r1) -- node[above right, font=\scriptsize] {N} (plano);
\draw[thick] (plano) -- node[above, font=\scriptsize] {1} (r2);
\draw[thick] (r2) -- node[above, font=\scriptsize] {N} (meta);
\draw[thick] (user) -- node[left, font=\scriptsize] {1} (r3);
\draw[thick] (r3) -- node[above, font=\scriptsize] {N} (sessao);
\draw[thick] (user) -- node[left, font=\scriptsize] {N} (r4);
\draw[thick] (r4) -- node[left, font=\scriptsize] {N} (mentoria);
\draw[thick] (user.east) -| (r5);
\draw[thick] (r5) -- node[right, font=\scriptsize] {N} (recurso);
\draw[thick] (user.south) |- (r6);
\draw[thick] (r6) -- node[above, font=\scriptsize] {N} (avaliacao);
\draw[thick] (user.south east) |- (r7);
\draw[thick] (r7) -- node[above, font=\scriptsize] {N} (badge);

\end{tikzpicture}
}%
\caption{Diagrama Entidade-Relacionamento}
\label{fig:diagrama-er}
\end{figure}

% ============================================================
% SEÇÃO 5 - TECNOLOGIAS
% ============================================================
\newpage
\section{Stack Tecnológica Recomendada}

\subsection{Visão Geral das Tecnologias}

\begin{center}
\begin{tabularx}{\textwidth}{|c|X|X|}
\hline
\rowcolor{faesaAzul}
\textcolor{white}{\textbf{Camada}} & \textcolor{white}{\textbf{Tecnologia}} & \textcolor{white}{\textbf{Justificativa}} \\
\hline
\textbf{Frontend} & React.js / Next.js 14+ & SSR/SSG para performance, TypeScript para segurança de tipos, ecossistema robusto \\
\hline
\textbf{Estilização} & Tailwind CSS + shadcn/ui & Design system consistente, utilitário, responsivo, dark mode nativo \\
\hline
\textbf{Backend} & Node.js + Express / NestJS & Mesmo ecossistema JavaScript, alta performance com V8, amplamente adotado \\
\hline
\textbf{Banco de Dados} & PostgreSQL 16+ & Relacional, robusto, suporte a JSON, extensível, open-source \\
\hline
\textbf{Cache} & Redis & Cache em memória para sessões, notificações em tempo real \\
\hline
\textbf{ORM} & Prisma & Type-safe, migrations automáticas, integração nativa com TypeScript \\
\hline
\textbf{Autenticação} & NextAuth.js / OAuth 2.0 & SSO com provedor FAESA, JWT, refresh tokens \\
\hline
\textbf{Real-time} & Socket.io & WebSockets para chat, notificações em tempo real, fórum ao vivo \\
\hline
\textbf{Deploy} & Vercel + Docker & CI/CD automatizado, preview deployments, escalabilidade \\
\hline
\textbf{Testes} & Jest + Cypress + Playwright & Unitários, integração e E2E para cobertura completa \\
\hline
\textbf{Monitoramento} & Sentry + Grafana & Error tracking, logs, métricas de performance \\
\hline
\end{tabularx}
\end{center}

\subsection{Diagrama de Tecnologias}

\begin{figure}[H]
\centering
\begin{tikzpicture}[
    tech/.style={draw, fill=#1!15, rounded corners=8pt, minimum width=2.5cm, 
                 minimum height=1cm, font=\small\bfseries, drop shadow, text centered},
    every node/.style={font=\small}
]
% Frontend
\node[tech=blue] (next) at (0, 0) {Next.js};
\node[tech=blue] (react) at (3.5, 0) {React};
\node[tech=cyan] (tw) at (7, 0) {Tailwind};
\node[tech=blue] (ts) at (10.5, 0) {TypeScript};

% Backend
\node[tech=green] (node) at (0, -2.5) {Node.js};
\node[tech=green] (nest) at (3.5, -2.5) {NestJS};
\node[tech=green] (prisma) at (7, -2.5) {Prisma};
\node[tech=green] (socket) at (10.5, -2.5) {Socket.io};

% Data
\node[tech=orange] (pg) at (0, -5) {PostgreSQL};
\node[tech=red] (redis) at (3.5, -5) {Redis};
\node[tech=orange] (s3) at (7, -5) {Cloud Storage};
\node[tech=purple] (sentry) at (10.5, -5) {Sentry};

% Infra
\node[tech=gray] (docker) at (2, -7.5) {Docker};
\node[tech=gray] (vercel) at (5.5, -7.5) {Vercel};
\node[tech=gray] (github) at (9, -7.5) {GitHub Actions};

% Labels
\node[font=\bfseries, text=blue!70] at (-3, 0) {Frontend};
\node[font=\bfseries, text=green!50!black] at (-3, -2.5) {Backend};
\node[font=\bfseries, text=orange!70!black] at (-3, -5) {Dados};
\node[font=\bfseries, text=gray!70!black] at (-3, -7.5) {Infra};

% Linhas horizontais
\draw[dashed, color=gray!50] (-4, -1.25) -- (12.5, -1.25);
\draw[dashed, color=gray!50] (-4, -3.75) -- (12.5, -3.75);
\draw[dashed, color=gray!50] (-4, -6.25) -- (12.5, -6.25);

\end{tikzpicture}
\caption{Mapa de Tecnologias por Camada}
\label{fig:tech-stack}
\end{figure}

% ============================================================
% SEÇÃO 6 - WIREFRAMES / PROTÓTIPOS
% ============================================================
\newpage
\section{Protótipos de Interface}

\subsection{Estrutura de Páginas}

A seguir, a representação esquemática das principais telas do site.

\subsubsection{Tela de Login}

\begin{figure}[H]
\centering
\begin{tikzpicture}
% Frame
\draw[thick, rounded corners=5pt] (0,0) rectangle (12, 8);

% Header
\fill[faesaAzul, rounded corners=5pt] (0, 7) rectangle (12, 8);
\node[text=white, font=\bfseries] at (6, 7.5) {FAESA --- Acolhimento Estudantil};

% Logo
\draw[fill=faesaAzul!10, rounded corners=3pt] (4, 5) rectangle (8, 6.5);
\node[font=\large\bfseries, text=faesaAzul] at (6, 5.75) {FAESA};

% Campos
\draw[fill=gray!10, rounded corners=2pt] (3, 4) rectangle (9, 4.5);
\node[font=\small, text=gray] at (6, 4.25) {E-mail institucional (@faesa.br)};

\draw[fill=gray!10, rounded corners=2pt] (3, 3) rectangle (9, 3.5);
\node[font=\small, text=gray] at (6, 3.25) {Senha};

% Botão
\draw[fill=faesaAzul, rounded corners=3pt] (4, 1.8) rectangle (8, 2.5);
\node[text=white, font=\small\bfseries] at (6, 2.15) {ENTRAR};

% Link
\node[font=\tiny, text=faesaAzulClaro] at (6, 1.3) {Primeiro acesso? Cadastre-se com sua matrícula};

% Rodapé
\draw[fill=faesaCinza!20] (0, 0) rectangle (12, 0.5);
\node[font=\tiny, text=faesaCinza] at (6, 0.25) {\textcopyright\ 2026 FAESA Centro Universitário};
\end{tikzpicture}
\caption{Wireframe --- Tela de Login}
\label{fig:wireframe-login}
\end{figure}

\subsubsection{Tela de Dashboard (Home)}

\begin{figure}[H]
\centering
\begin{tikzpicture}[scale=0.9, every node/.style={scale=0.9}]
% Frame
\draw[thick, rounded corners=5pt] (0,0) rectangle (16, 11);

% Header
\fill[faesaAzul, rounded corners=5pt] (0, 10) rectangle (16, 11);
\node[text=white, font=\bfseries] at (3, 10.5) {\faIcon{home} Dashboard};
\node[text=white, font=\small] at (13, 10.5) {\faIcon{bell} \faIcon{user-circle} Lucas};

% Sidebar
\fill[faesaAzul!10] (0, 0) rectangle (3.5, 10);
\node[font=\small\bfseries, text=faesaAzul] at (1.75, 9.5) {Menu};
\node[font=\scriptsize, text=faesaAzul] at (1.75, 8.8) {\faIcon{tachometer-alt} Dashboard};
\node[font=\scriptsize, text=faesaCinza] at (1.75, 8.2) {\faIcon{book} Plano de Estudos};
\node[font=\scriptsize, text=faesaCinza] at (1.75, 7.6) {\faIcon{brain} Concentração};
\node[font=\scriptsize, text=faesaCinza] at (1.75, 7.0) {\faIcon{folder-open} Biblioteca};
\node[font=\scriptsize, text=faesaCinza] at (1.75, 6.4) {\faIcon{users} Mentoria};
\node[font=\scriptsize, text=faesaCinza] at (1.75, 5.8) {\faIcon{comments} Fórum};
\node[font=\scriptsize, text=faesaCinza] at (1.75, 5.2) {\faIcon{trophy} Badges};
\node[font=\scriptsize, text=faesaCinza] at (1.75, 4.6) {\faIcon{heart} Bem-estar};

% Cards de progresso
\draw[fill=faesaVerde!10, rounded corners=3pt] (4, 8) rectangle (7.5, 9.8);
\node[font=\scriptsize\bfseries, text=faesaVerde] at (5.75, 9.4) {Metas Concluídas};
\node[font=\Large\bfseries, text=faesaVerde] at (5.75, 8.7) {12/20};

\draw[fill=faesaAzulClaro!10, rounded corners=3pt] (8, 8) rectangle (11.5, 9.8);
\node[font=\scriptsize\bfseries, text=faesaAzulClaro] at (9.75, 9.4) {Horas de Estudo};
\node[font=\Large\bfseries, text=faesaAzulClaro] at (9.75, 8.7) {45h};

\draw[fill=faesaLaranja!10, rounded corners=3pt] (12, 8) rectangle (15.5, 9.8);
\node[font=\scriptsize\bfseries, text=faesaLaranja] at (13.75, 9.4) {Sessões de Foco};
\node[font=\Large\bfseries, text=faesaLaranja] at (13.75, 8.7) {28};

% Gráfico (simulado)
\draw[fill=gray!5, rounded corners=3pt] (4, 3.5) rectangle (10, 7.5);
\node[font=\scriptsize\bfseries, text=faesaAzul] at (7, 7.2) {Progresso Semanal};
% Barras do gráfico
\foreach \x/\h/\label in {4.5/2.5/Seg, 5.3/3.0/Ter, 6.1/1.8/Qua, 6.9/3.5/Qui, 7.7/2.8/Sex, 8.5/1.2/Sáb, 9.3/0.5/Dom} {
    \fill[faesaAzul!60, rounded corners=1pt] (\x, 3.8) rectangle (\x+0.5, 3.8+\h);
    \node[font=\tiny, text=faesaCinza] at (\x+0.25, 3.5) {\label};
}

% Próximas atividades
\draw[fill=gray!5, rounded corners=3pt] (10.5, 3.5) rectangle (15.5, 7.5);
\node[font=\scriptsize\bfseries, text=faesaAzul] at (13, 7.2) {Próximas Atividades};
\draw[fill=faesaAzul!5, rounded corners=2pt] (10.8, 6.0) rectangle (15.2, 6.8);
\node[font=\tiny] at (13, 6.4) {\faIcon{clock} Sessão Pomodoro --- 14:00};
\draw[fill=faesaLaranja!5, rounded corners=2pt] (10.8, 5.0) rectangle (15.2, 5.8);
\node[font=\tiny] at (13, 5.4) {\faIcon{book} Estudar Cálculo --- 16:00};
\draw[fill=faesaVerde!5, rounded corners=2pt] (10.8, 4.0) rectangle (15.2, 4.8);
\node[font=\tiny] at (13, 4.4) {\faIcon{users} Mentoria --- Amanhã 10:00};

% Badges recentes
\draw[fill=gray!5, rounded corners=3pt] (4, 0.5) rectangle (15.5, 3);
\node[font=\scriptsize\bfseries, text=faesaAzul] at (5.5, 2.7) {Badges Recentes};
\foreach \x/\cor/\nome in {5.5/faesaVerde/Focado, 8/faesaAzulClaro/Dedicado, 10.5/faesaLaranja/Mentor, 13/purple/Explorador} {
    \node[draw=\cor, fill=\cor!20, circle, minimum size=1.2cm, font=\tiny\bfseries, text=\cor] at (\x, 1.5) {\nome};
}

\end{tikzpicture}
\caption{Wireframe --- Dashboard Principal}
\label{fig:wireframe-dashboard}
\end{figure}

% ============================================================
% SEÇÃO 7 - MÓDULOS DETALHADOS
% ============================================================
\newpage
\section{Detalhamento dos Módulos}

\subsection{Módulo 1 --- Plano de Estudos}

\begin{destaque}[Funcionalidades do Plano de Estudos]
O módulo de Plano de Estudos permite ao estudante:
\begin{itemize}
    \item Criar planos semanais ou mensais vinculados às suas disciplinas;
    \item Definir metas SMART (Específicas, Mensuráveis, Alcançáveis, Relevantes, Temporais);
    \item Utilizar calendário interativo com drag-and-drop;
    \item Receber sugestões automáticas baseadas no curso e período;
    \item Acompanhar o percentual de cumprimento dos planos;
    \item Exportar o plano em formato PDF ou integrar com Google Calendar.
\end{itemize}
\end{destaque}

\subsubsection{Fluxo de Criação de Plano}

\begin{figure}[H]
\centering
\begin{tikzpicture}[
    step/.style={draw=faesaAzul, fill=faesaAzul!10, rounded corners=5pt, 
                 minimum width=2.5cm, minimum height=1.2cm, text centered, 
                 font=\scriptsize, text width=2.3cm},
    arrow/.style={-{Stealth[length=3mm]}, thick, color=faesaAzul}
]

\node[step] (s1) at (0, 0) {Selecionar Disciplinas do Período};
\node[step] (s2) at (3.5, 0) {Definir Horários Disponíveis};
\node[step] (s3) at (7, 0) {Criar Metas Semanais};
\node[step] (s4) at (10.5, 0) {Configurar Lembretes};
\node[step] (s5) at (14, 0) {Ativar Plano};

\draw[arrow] (s1) -- (s2);
\draw[arrow] (s2) -- (s3);
\draw[arrow] (s3) -- (s4);
\draw[arrow] (s4) -- (s5);

% Números
\foreach \x/\n in {0/1, 3.5/2, 7/3, 10.5/4, 14/5} {
    \node[circle, fill=faesaAzul, text=white, font=\tiny\bfseries, minimum size=0.5cm] at (\x, 0.9) {\n};
}

\end{tikzpicture}
\caption{Fluxo de Criação de Plano de Estudos}
\label{fig:fluxo-plano}
\end{figure}

\subsection{Módulo 2 --- Exercícios de Concentração}

\begin{destaque}[Tipos de Exercícios Disponíveis]
\begin{enumerate}
    \item \textbf{Técnica Pomodoro} --- Ciclos de 25 min de foco + 5 min de pausa, com timer integrado;
    \item \textbf{Mindfulness} --- Meditações guiadas de 5, 10 ou 15 minutos com áudio;
    \item \textbf{Exercícios Respiratórios} --- Técnica 4-7-8 com animação visual;
    \item \textbf{White Noise / Sons Ambientes} --- Biblioteca de sons para concentração (chuva, floresta, café);
    \item \textbf{Focus Mode} --- Bloqueio de notificações e gamificação do tempo de foco.
\end{enumerate}
\end{destaque}

\subsection{Módulo 3 --- Mentoria}

O sistema de mentoria conecta estudantes veteranos a calouros, promovendo troca de experiências e suporte acadêmico.

\begin{center}
\begin{tabularx}{\textwidth}{|l|X|}
\hline
\rowcolor{faesaAzul}
\textcolor{white}{\textbf{Aspecto}} & \textcolor{white}{\textbf{Detalhamento}} \\
\hline
\textbf{Quem pode ser mentor?} & Estudantes a partir do 5º período com CRA $\geq$ 7,0 \\
\hline
\textbf{Como funciona?} & O calouro solicita mentoria; o sistema faz match por curso e disponibilidade \\
\hline
\textbf{Frequência} & Mínimo de 2 encontros mensais (presenciais ou online) \\
\hline
\textbf{Avaliação} & Ambas as partes avaliam ao final de cada ciclo (0-10) \\
\hline
\textbf{Benefício mentor} & Pontos de gamificação, certificado de horas complementares \\
\hline
\end{tabularx}
\end{center}

\subsection{Módulo 4 --- Estudos Fora da FAESA}

\begin{destaque}[Recursos Externos Recomendados]
O módulo cataloga e recomenda atividades de estudo externas:
\begin{itemize}
    \item \textbf{Cursos Online:} Coursera, edX, Alura, Rocketseat, Khan Academy;
    \item \textbf{Certificações:} AWS, Google, Microsoft, Oracle;
    \item \textbf{Eventos:} Hackathons, meetups, conferências da área;
    \item \textbf{Leitura:} Artigos científicos, blogs técnicos, livros recomendados;
    \item \textbf{Projetos:} Desafios no GitHub, contribuições open source.
\end{itemize}
Cada recurso possui avaliação da comunidade e tag de relevância por curso.
\end{destaque}

\subsection{Módulo 5 --- Gamificação}

O sistema de gamificação incentiva o engajamento contínuo dos estudantes:

\begin{center}
\begin{tabular}{|l|l|l|}
\hline
\rowcolor{faesaAzul}
\textcolor{white}{\textbf{Badge}} & \textcolor{white}{\textbf{Critério}} & \textcolor{white}{\textbf{Pontos}} \\
\hline
\faIcon{star} Primeiro Passo & Completar o cadastro e primeiro plano & 50 pts \\
\hline
\faIcon{fire} Focado & 7 dias consecutivos de sessões Pomodoro & 100 pts \\
\hline
\faIcon{graduation-cap} Dedicado & Cumprir 100\% das metas semanais & 150 pts \\
\hline
\faIcon{users} Mentor & Completar 1 ciclo de mentoria & 200 pts \\
\hline
\faIcon{book-reader} Explorador & Acessar 10 recursos da biblioteca & 75 pts \\
\hline
\faIcon{brain} Zen Master & 30 sessões de mindfulness & 120 pts \\
\hline
\faIcon{trophy} Top Estudante & Acumular 1000 pontos totais & 300 pts \\
\hline
\end{tabular}
\end{center}

% ============================================================
% SEÇÃO 8 - CRONOGRAMA
% ============================================================
\newpage
\section{Cronograma de Desenvolvimento}

O desenvolvimento do site está dividido em sprints de 2 semanas, seguindo a metodologia ágil Scrum.

\begin{figure}[H]
\centering
\resizebox{\textwidth}{!}{%
\begin{tikzpicture}[
    sprint/.style={draw=faesaAzul, fill=#1, rounded corners=3pt, 
                   minimum height=0.7cm, font=\scriptsize, text=black, anchor=west},
    label/.style={font=\scriptsize\bfseries, text=faesaAzul, anchor=east}
]

% Eixo de tempo
\draw[thick, -{Stealth}] (0, 0) -- (16, 0);
\foreach \x/\m in {0/Mar, 2/Abr, 4/Mai, 6/Jun, 8/Jul, 10/Ago, 12/Set, 14/Out} {
    \draw[thick] (\x, -0.2) -- (\x, 0.2);
    \node[font=\scriptsize, below] at (\x, -0.3) {\m};
}

% Sprints
\node[label] at (-0.2, 5) {Planejamento};
\node[sprint=faesaAzul!30, minimum width=3cm] at (0, 5) {Requisitos e Design};

\node[label] at (-0.2, 4) {Sprint 1-2};
\node[sprint=faesaVerde!30, minimum width=4cm] at (2, 4) {Auth + Cadastro + Dashboard};

\node[label] at (-0.2, 3) {Sprint 3-4};
\node[sprint=faesaVerde!30, minimum width=4cm] at (4, 3) {Plano de Estudos + Cronograma};

\node[label] at (-0.2, 2) {Sprint 5-6};
\node[sprint=faesaLaranja!30, minimum width=4cm] at (6, 2) {Concentração + Gamificação};

\node[label] at (-0.2, 1) {Sprint 7-8};
\node[sprint=faesaLaranja!30, minimum width=4cm] at (8, 1) {Mentoria + Fórum + Chat};

\node[label] at (-0.2, -1) {Sprint 9-10};
\node[sprint=purple!20, minimum width=4cm] at (10, -1) {Biblioteca + Atividades Ext.};

\node[label] at (-0.2, -2) {Sprint 11-12};
\node[sprint=red!20, minimum width=3cm] at (12, -2) {Testes + Deploy + Docs};

% Marco
\node[draw=faesaVerde, fill=faesaVerde, star, star points=5, minimum size=0.4cm] at (15, -2) {};
\node[font=\scriptsize\bfseries, text=faesaVerde, right] at (15.3, -2) {Go-Live!};

\end{tikzpicture}
}%
\caption{Cronograma de Desenvolvimento (Gantt Simplificado)}
\label{fig:cronograma}
\end{figure}

\begin{center}
\begin{tabularx}{\textwidth}{|c|c|X|}
\hline
\rowcolor{faesaAzul}
\textcolor{white}{\textbf{Fase}} & \textcolor{white}{\textbf{Período}} & \textcolor{white}{\textbf{Entregas}} \\
\hline
Planejamento & Mar 2026 & Requisitos, protótipos, arquitetura, setup do ambiente \\
\hline
Sprint 1--2 & Abr 2026 & Autenticação SSO, cadastro, layout base, dashboard \\
\hline
Sprint 3--4 & Mai 2026 & CRUD plano de estudos, calendário interativo, metas \\
\hline
Sprint 5--6 & Jun 2026 & Timer Pomodoro, mindfulness, gamificação base \\
\hline
Sprint 7--8 & Jul 2026 & Sistema de mentoria, fórum, chat em tempo real \\
\hline
Sprint 9--10 & Ago 2026 & Biblioteca de recursos, atividades extracurriculares \\
\hline
Sprint 11--12 & Set--Out 2026 & Testes E2E, otimização de performance, deploy final \\
\hline
\end{tabularx}
\end{center}

% ============================================================
% SEÇÃO 9 - CONCLUSÃO
% ============================================================
\newpage
\section{Conclusão}

O \textbf{Site de Acolhimento da FAESA} é uma plataforma estratégica para o sucesso acadêmico dos estudantes, integrando em um único ambiente digital as ferramentas necessárias para:

\begin{itemize}[leftmargin=2cm]
    \item[\faIcon{check-circle}] \textbf{Organização acadêmica} por meio de planos de estudo personalizados;
    \item[\faIcon{check-circle}] \textbf{Desenvolvimento pessoal} com exercícios de concentração e mindfulness;
    \item[\faIcon{check-circle}] \textbf{Conexão social} através de mentoria e fórum de discussão;
    \item[\faIcon{check-circle}] \textbf{Aprendizado contínuo} com biblioteca de recursos e atividades externas;
    \item[\faIcon{check-circle}] \textbf{Motivação} por meio de gamificação e acompanhamento de progresso;
    \item[\faIcon{check-circle}] \textbf{Bem-estar} com avaliações periódicas e suporte psicopedagógico.
\end{itemize}

O projeto segue boas práticas de engenharia de software, com arquitetura modular, stack tecnológica moderna e desenvolvimento ágil, garantindo que a plataforma possa evoluir continuamente para atender às necessidades dos mais de 75 mil alunos formados pela instituição e dos futuros estudantes que ingressarão na FAESA.

\begin{destaque}[Próximos Passos]
\begin{enumerate}
    \item Validação dos requisitos com stakeholders (coordenadores e estudantes);
    \item Criação de protótipos navegáveis no Figma;
    \item Setup do ambiente de desenvolvimento e CI/CD;
    \item Início da Sprint 1 com a funcionalidade de autenticação;
    \item Testes de usabilidade com grupo piloto de 50 estudantes.
\end{enumerate}
\end{destaque}

% ============================================================
% REFERÊNCIAS
% ============================================================
\newpage
\section{Referências}

\begin{enumerate}[leftmargin=2cm]
    \item FAESA Centro Universitário. \textit{Sobre Nós}. Disponível em: \url{https://www.faesa.br/sobre-nos/}. Acesso em: fev. 2026.
    
    \item SOMMERVILLE, Ian. \textit{Engenharia de Software}. 10ª ed. São Paulo: Pearson, 2018.
    
    \item PRESSMAN, Roger S. \textit{Engenharia de Software: Uma Abordagem Profissional}. 8ª ed. Porto Alegre: AMGH, 2016.
    
    \item W3C. \textit{Web Content Accessibility Guidelines (WCAG) 2.1}. Disponível em: \url{https://www.w3.org/TR/WCAG21/}. Acesso em: fev. 2026.
    
    \item REACT. \textit{React Documentation}. Disponível em: \url{https://react.dev/}. Acesso em: fev. 2026.
    
    \item NEXT.JS. \textit{Next.js Documentation}. Disponível em: \url{https://nextjs.org/docs}. Acesso em: fev. 2026.
    
    \item CIRILLO, Francesco. \textit{The Pomodoro Technique}. Berlin: FC Garage, 2006.
    
    \item BRASIL. \textit{Lei nº 13.709/2018 --- Lei Geral de Proteção de Dados Pessoais (LGPD)}. Diário Oficial da União, 2018.
    
    \item FOWLER, Martin. \textit{UML Distilled: A Brief Guide to the Standard Object Modeling Language}. 3ª ed. Addison-Wesley, 2003.
    
    \item NIELSEN, Jakob. \textit{Usability Engineering}. San Francisco: Morgan Kaufmann, 1993.
\end{enumerate}

% ============================================================
\end{document}
